\section{Fractals: Mandelbrot and Julia Sets}

\begin{frame}{What Is a Fractal?}
  \begin{itemize}
    \item A \textbf{fractal} is a geometric object that exhibits \textbf{self-similarity} at different scales.
    \item Zooming in reveals structure that resembles the whole.
    \item Fractals often arise from simple iterative rules applied repeatedly.
    \item They are intimately connected to chaos: the boundary of the Mandelbrot set, for example, encodes all possible behaviours of a simple quadratic iteration.
    \pause
    \item Examples in nature:
    \begin{itemize}
      \item Coastlines, mountain ranges
      \item Branching of trees, blood vessels, lungs
      \item Snowflakes, lightning bolts
      \item Romanesco broccoli
    \end{itemize}
  \end{itemize}
\end{frame}

\begin{frame}{Iteration in the Complex Plane}
  \begin{itemize}
    \item Recall: a complex number $c = x + iy$ is a point in the 2D plane.
    \item We study the iteration:
    \[
      z_{n+1} = z_n^2 + c
    \]
    \item Starting from $z_0 = 0$, does the sequence $\{z_n\}$ stay bounded or escape to infinity?
    \item If $|z_n| > 2$ at any point, the sequence will diverge.
    \pause
    \item Two key objects arise from this single equation:
    \begin{itemize}
      \item \textbf{Mandelbrot set}: fix $z_0 = 0$, vary $c$. Colour by whether $c$ causes divergence.
      \item \textbf{Julia set}: fix $c$, vary $z_0$. Colour by whether $z_0$ causes divergence.
    \end{itemize}
  \end{itemize}
\end{frame}

\begin{frame}{The Mandelbrot Set}
  \begin{itemize}
    \item The Mandelbrot set is the set of all complex numbers $c$ for which the iteration $z_{n+1} = z_n^2 + c$ (starting from $z_0 = 0$) does \textbf{not} diverge.
    \item Points \textbf{inside} the set are coloured black.
    \item Points \textbf{outside} are coloured by how quickly they escape (the \textbf{escape time}).
    \pause
    \item Remarkable properties:
    \begin{itemize}
      \item The boundary has \textbf{infinite complexity} --- zooming in reveals ever more detail.
      \item Small copies of the whole set appear at all scales.
      \item The area is finite, but the boundary has infinite length.
      \item The boundary has a fractal dimension $\approx 2$.
    \end{itemize}
  \end{itemize}
\end{frame}

\begin{frame}{Julia Sets}
  \begin{itemize}
    \item For a fixed $c$, the \textbf{Julia set} is the boundary between initial values $z_0$ that diverge and those that stay bounded.
    \item Each value of $c$ produces a different Julia set.
    \pause
    \item There is a deep connection between Mandelbrot and Julia sets:
    \begin{itemize}
      \item If $c$ is \textbf{inside} the Mandelbrot set, the Julia set is a connected, often beautiful fractal.
      \item If $c$ is \textbf{outside} the Mandelbrot set, the Julia set is a disconnected dust of points (a \textbf{Cantor set}).
    \end{itemize}
    \item The Mandelbrot set is therefore a ``map'' of all possible Julia sets --- it tells us which values of $c$ produce connected fractals.
  \end{itemize}
\end{frame}

\begin{frame}{Connection to Chaos}
  \begin{itemize}
    \item The Mandelbrot set and Julia sets arise from the \textbf{same iteration} that produces chaos in the logistic map.
    \item In fact, the logistic map $X_{N+1} = r\,X_N(1 - X_N)$ is related to $z_{n+1} = z_n^2 + c$ by a change of variables.
    \pause
    \item The bifurcation diagram of the logistic map corresponds to a slice through the Mandelbrot set along the real axis.
    \item The Feigenbaum constant $\delta \approx 4.669$ also governs the scaling of structures in the Mandelbrot set.
    \item All three topics of this lecture --- the logistic map, the Lorenz system, and fractals --- are manifestations of the same underlying mathematics of \textbf{nonlinear dynamics}.
  \end{itemize}
\end{frame}

\begin{frame}{Summary: Chaos in Three Forms}
  \begin{center}
    \begin{tabular}{lll}
      \textbf{System} & \textbf{Type} & \textbf{Key feature} \\
      \hline
      Logistic map & Discrete, 1D & Period doubling to chaos \\
      Lorenz system & Continuous, 3D & Strange attractor \\
      Mandelbrot set & Complex iteration & Fractal boundary \\
    \end{tabular}
  \end{center}

  \vfill
  All three share:
  \begin{itemize}
    \item Sensitivity to initial conditions
    \item Deterministic yet unpredictable behaviour
    \item Rich structure from simple rules
  \end{itemize}
\end{frame}
